% =============================================================================
% Chapter 1: Introduction
% =============================================================================

\chapter{Introduction}
\label{ch:introduction}

\section{Purpose of this Document}

This Software Requirements Specification (SRS) defines the requirements,
scope, architecture, and contribution standards for \itc{}
(\code{import itaca as itc})~\citep{python}, a Python library for rigorous
engineering data management, analysis, and computation, with a primary focus
on aerospace applications.

The document is the authoritative reference for:
\begin{itemize}[noitemsep,topsep=2pt]
  \item what \itc{} must do (functional requirements);
  \item what \itc{} must not do or be (non-requirements and out-of-scope items);
  \item how \itc{} must behave (non-functional requirements);
  \item how the codebase is structured (architecture and class hierarchy);
  \item how contributors must work (coding standards, testing, versioning).
\end{itemize}

\begin{noteinfo}
This SRS is a living document. As use cases consolidate, requirements are
added, refined, or deprecated. Every change increments the document version
and is recorded in the revision history table. Deprecated requirements are
never deleted: they are kept in place and tagged \deprecated{}.
\end{noteinfo}

\begin{noteinfo}[Canonical baseline]
This version is the first tracked in the research workspace and supersedes
the founding project instructions of May 2026. Where the founding material
and this document diverge, this document prevails. Canonical terms fixed
here: the equation file extension is \code{.itceq}; the string derivation
method is \code{db.compute}; \code{Processor} is a typing protocol;
correlated uncertainty ships in the first release; the figure wrapper is
\code{ItcFigure}. The founding material is preserved unchanged as a
historical record in the project archive.
\end{noteinfo}

\begin{noteinfo}[On implementation tooling]
\itc{} is being implemented with assistance from large language model tooling,
notably Claude Code. Every implementation derived from these requirements goes
through human double-check by the project author before being accepted. The
SRS describes the desired behavior; the code is verified against the SRS, and
the SRS is corrected only when a requirement itself is found to be wrong or
ambiguous.
\end{noteinfo}

\section{Context and Motivation}

Experimental and numerical aerospace data (wind tunnel campaigns, propeller performance tests, computational fluid dynamics post-processing, in-flight testing, and conceptual aircraft analyses) share a recurring set of
challenges that no existing general-purpose tool addresses in a unified,
rigorous way:

\begin{itemize}
  \item \textbf{Data is naturally multidimensional.} Results depend on angle
    of attack, sideslip angle, Mach number, Reynolds number, RPM, blade
    pitch, control surface deflections, radial position, azimuthal angle,
    and chord-wise position simultaneously.

  \item \textbf{Campaigns are incomplete by nature.} Not all combinations of
    dimensions are tested or simulated; missing data must be handled
    explicitly, not silently.

  \item \textbf{Results feed engineering decisions.} The pipeline from raw
    measurements to published coefficients to certification-relevant numbers
    must be fully traceable and reproducible. Traceability is not a feature;
    it is a prerequisite.

  \item \textbf{Uncertainty quantification is mandatory} for serious
    experimental work, not optional. It must propagate automatically through
    every derived quantity following the principles of the
    GUM~\citep{gum2008}, including covariance terms when input variables are
    correlated.

  \item \textbf{Workflows repeat across campaigns} with minor variations.
    Load, correct, derive coefficients, smooth, plot, export: the same
    skeleton is rewritten ad hoc by every team. It must instead be
    encapsulated and reusable through versioned equation files and named
    processors.

  \item \textbf{Engineering computations and experimental data must share a
    common substrate.} Performance analyses, stability and control sizing,
    and Blade Element Momentum Theory results must be generated in the same
    data structure used by experimental measurements, so that the same
    plotting, comparison, surrogate-modeling, and uncertainty machinery
    applies seamlessly to both.
\end{itemize}

\itc{} was conceived to address all of these challenges in a single, coherent
framework, designed from first principles around the needs of aerospace
experimental and numerical data analysis. While its design center is
aerospace, the core data structures and operations are general enough to
support other engineering and scientific domains where multidimensional
datasets, traceability, and uncertainty quantification are fundamental.

\section{Document Conventions}

\subsection{Requirement Identifiers}

Functional and non-functional requirements are identified as \reqid{REQ-XX},
with \code{XX} a zero-padded two-digit number. Non-requirements (explicit
out-of-scope statements) are identified as \nreqid{NREQ-XX}. Use cases are
\code{UC-XX} and design decisions are \ddid{DD-XX}. Identifiers are stable: a
requirement deprecated in a later version keeps its number forever.

Requirements and decisions are color-coded:
\begin{itemize}[noitemsep,topsep=2pt]
  \item \textcolor{itacaTangerine}{\textbf{Tangerine}}: functional or
        non-functional requirement (what the system shall do).
  \item \textcolor{itacaDeepRed}{\textbf{Deep red}}: non-requirement (what
        the system shall not do or be).
  \item \textcolor{itacaOrange}{\textbf{Orange}}: design decision (how a
        requirement is satisfied).
\end{itemize}

\subsection{Status Tags}

Each requirement carries one of the following status tags:

\begin{tabularx}{\textwidth}{@{}lX@{}}
\toprule
\textbf{Tag} & \textbf{Meaning}\\
\midrule
\stable      & Requirement is frozen; changes require formal review.\\
\draft       & Requirement is still being refined.\\
\implemented & Corresponding code exists, tested, and released.\\
\pending     & Not yet implemented.\\
\deprecated  & Superseded but kept for traceability.\\
\bottomrule
\end{tabularx}

\section{Definitions, Acronyms, and Abbreviations}

See the Acronyms and Glossary sections at the front of this document.

\section{References and Standards}

This document is informed by the following standards and references:

\begin{itemize}
  \item BIPM GUM (JCGM~100:2008)~\citep{gum2008}: normative reference for
    uncertainty propagation in \itc{}.
  \item BIPM GUM Supplement~1 (JCGM~101:2008)~\citep{gum2008s1}: Monte
    Carlo method, used in \itc{} for discrete-branch model handling.
  \item AIAA S-071A-1999 \emph{Assessment of Wind Tunnel Data
    Uncertainty}~\citep{aiaaS071A}: domain-specific application of the GUM
    for wind tunnel work.
  \item AIAA R-004A-1992 \emph{Recommended Practice for Atmospheric and Space
    Flight Vehicle Coordinate Systems}~\citep{aiaaR004A}: normative
    reference for axes and sign conventions.
  \item SAE AIR~4065 \emph{Propeller/Propfan In-Flight Thrust
    Determination}~\citep{saeAIR4065}: reference for the
    \code{FT\_thrust} processor.
  \item Barlow, Rae \& Pope, \emph{Low-Speed Wind Tunnel Testing},
    3rd~ed.~\citep{barlowRaePope}: reference for wind tunnel correction
    processors.
  \item Adkins \& Liebeck, \emph{Design of Optimum
    Propellers}~\citep{adkinsLiebeck}: reference for the
    \code{itc.aerospace.propeller} BEMT implementation.
  \item Python PEP~8 \emph{Style Guide for Python Code}~\citep{pep8}.
  \item NumPy~\citep{numpy}: the only mandatory scientific dependency in
    the \itc{} core.
  \item SMT: Surrogate Modelling Toolbox~\citep{smt}: backend for
    \code{itc.surrogate}.
  \item W3C PROV-DM \emph{Provenance Data Model}~\citep{provdm}: conceptual reference for the Provenance / History split, with native
    PROV-N and PROV-JSON export support in \itc{}.
\end{itemize}
